
\chapter[Phase-Derived Entropy Candidate]{Phase-Derived Entropy\\Candidate}
\section{Double-cover geometry}
The phase construction uses a full phase coordinate $\Theta\in[0,4\pi)$ and its projected phase
\begin{equation}
\theta=\Theta\bmod 2\pi.
\end{equation}
The hidden sheet label is
\begin{equation}
b=\left\lfloor\frac{\Theta\bmod4\pi}{2\pi}\right\rfloor\in\{0,1\}.
\end{equation}
For a balanced two-lift conditional distribution,
\begin{equation}
H(b)=-\sum_{b\in\{0,1\}}\frac12\ln\frac12=\ln2,
\end{equation}
and equivalently
\begin{equation}
\ln\frac{4\pi}{2\pi}=\ln2.
\end{equation}
The implementation keeps this geometric multiplicity distinct from empirical cryptographic unpredictability. Cryptographic admission is therefore based on measured min-entropy, lag prediction, run-length and adaptive-proportion health tests.

\section{Live evidence}
A recorded live run sampled 48 unique 36D frames and produced 1702 raw noise bits. The measured health snapshot was
\begin{align}
\Hmin &\approx 0.899784\ \text{bit/sample},\\
H_{\mathrm{Markov}} &\approx 0.926172,\\
H_{\mathrm{lag}} &\approx 0.899784,\\
\text{bias} &\approx 0.007638,\\
A_{\mathrm{pred}} &\approx 0.035967.
\end{align}
The estimated aggregate min-entropy budget was 1531.43 bits. Repetition-count and adaptive-proportion health gates passed. The source class recorded in provenance is \code{ENDOGENOUS\_RUNTIME\_DYNAMICS}.

\begin{resultbox}[Entropy accounting rule]
The measured entropy budget is assigned to the conditioned seed input. SHAKE256 provides deterministic, purpose-separated expansion. Expansion length does not increase the measured entropy budget.
\end{resultbox}

\chapter{36D Entropy Topology}
\section{Per-coordinate measurements}
For coordinate streams $X_i$, $i=0,\ldots,35$, the topology stage measures
\begin{equation}
\Hmin^{(i)},\qquad I(X_i;X_j),\qquad A_{\mathrm{pred}}(i,j).
\end{equation}
The selector constructs a dynamic subset $S\subset\{0,\ldots,35\}$ under thresholds on minimum entropy, pairwise mutual information, and cross-coordinate prediction advantage.

A graph representation is
\begin{equation}
G_E=(V,E,w),\qquad w_{ij}=I(X_i;X_j),
\end{equation}
with optional predictive edge metadata. The selection mask is recomputed per observation window.

\section{Recorded topology windows}
One 72-frame live window selected 13 of 36 coordinates with a selected budget of 592.73 bits. A subsequent validation window selected 17 coordinates. For that selected set,
\begin{align}
\max I(X_i;X_j)&\approx0.03832\ \text{bit},\\
\max A_{\mathrm{pred}}(i,j)&\approx0.11429,\\
\Hmin^{\mathrm{combined}}&\approx0.89973\ \text{bit/sample},\\
N_{\mathrm{combined}}&=1190.
\end{align}
The combined health gate passed.

\begin{figure}[ht]
\centering
\begin{tikzpicture}[scale=.9,every node/.style={font=\sffamily\scriptsize}]
\foreach \i in {0,...,11}{\node[circle,draw=DNCblue,fill=DNClight,minimum size=7mm] (n\i) at ({3*cos(30*\i)},{3*sin(30*\i)}) {\i};}
\foreach \a/\b in {0/1,1/2,2/4,4/5,5/7,7/8,8/10,10/11,11/0,1/7,3/9,2/8}{\draw[DNCcyan!70,line width=.8pt] (n\a)--(n\b);}
\node[align=center,font=\sffamily\small] at (0,0) {illustrative\\entropy-dependence\\graph};
\end{tikzpicture}
\caption{Graph view of inter-coordinate information dependence. The production implementation evaluates all coordinate pairs; this drawing is schematic.}
\end{figure}

\chapter[Controlled Cryptanalysis Challenge v0.1]{Controlled Cryptanalysis\\Challenge v0.1}
\section{Experimental question}
The challenge asks whether the adaptive layer reduces adversarial risk on unseen ciphertexts when compared with a frozen static DNCrypt anchor under an equal attack budget. The frozen anchor uses the same standard AEAD/KDF suite in both conditions; the experimental variable is adaptive genome selection.

\section{Train/holdout protocol}
Each replicate uses 64 attack-train messages and 192 independent holdout messages, each 256 bytes. The adaptive condition receives three generations with 24 candidates per generation. Candidate selection may inspect attack-train metrics only. The primary endpoint is independent holdout adversarial risk.

The frozen configuration is summarized in Table~\ref{tab:challenge}.
\begin{table}[ht]
\centering
\caption{Cryptanalysis Challenge v0.1 configuration.}\label{tab:challenge}
\begin{tabular}{lr}
\toprule
Parameter & Value\\
\midrule
Replicates & 12\\
Generations per replicate & 3\\
Population size & 24\\
Train messages & 64\\
Holdout messages & 192\\
Message size & 256 bytes\\
Minimum throughput ratio & 0.5\\
Primary endpoint & holdout adversarial risk\\
\bottomrule
\end{tabular}
\end{table}

\section{Reference result}
Across 12 frozen replicates, adaptive DNCrypt won 4 holdout comparisons and the static anchor won 8. The mean holdout risk delta was
\begin{equation}